% >>> typage-template cvpr@2026.10.1 variant=review
\documentclass[10pt,twocolumn,letterpaper]{article}
\usepackage[review]{cvpr}
% cvpr selects Times (ptm); under XeTeX that needs the T1 encoding.
\usepackage[T1]{fontenc}
\def\confName{CVPR}
\def\confYear{2027}
\bibliographystyle{ieeenat_fullname}
% <<< typage-template
\usepackage{booktabs}
% hyperref is strongly recommended by the author kit, loaded last.
\definecolor{cvprblue}{rgb}{0.21,0.49,0.74}
\usepackage[pagebackref,breaklinks,colorlinks,allcolors=cvprblue]{hyperref}

\title{Paper Title}

\author{First Author\\
Institution\\
City, Country\\
{\tt\small first.author@example.com}
\and
Second Author\\
Institution\\
City, Country\\
{\tt\small second.author@example.com}
}

\begin{document}
\maketitle

\begin{abstract}
State the problem, the approach and the main result in a single
paragraph.
\end{abstract}

\section{Introduction}
\label{sec:intro}
Introduce the problem and list the contributions. Cite prior work by
number~\cite{lamport1994latex,knuth1984texbook}.

\section{Method}
\label{sec:method}
Describe the method. Number the equations you refer to:
\begin{equation}
  \label{eq:objective}
  \min_{\theta} \; \frac{1}{n} \sum_{i=1}^{n} \ell\bigl(f_\theta(x_i), y_i\bigr) .
\end{equation}

\section{Experiments}
\label{sec:experiments}
Table~\ref{tab:results} reports the results of
optimizing~\eqref{eq:objective}. Describe the data, the baselines and the
compute used, and report the variation across runs.

\begin{table}[!htb]
  \caption{Results on the example benchmark (mean $\pm$ standard deviation over five runs).}
  \label{tab:results}
  \centering
  \begin{tabular}{lcc}
    \toprule
    Method   & Accuracy (\%)  & Time (s) \\
    \midrule
    Baseline & 81.2 $\pm$ 0.4 & 12.4     \\
    Ours     & 84.7 $\pm$ 0.3 & 10.9     \\
    \bottomrule
  \end{tabular}
\end{table}

\section{Conclusion}
Conclude and outline future work.

{\small
\bibliography{refs}
}

\end{document}
